\chapter{Spécification des besoins}

\section*{Introduction}
Après avoir présenté le contexte et les objectifs du projet dans le chapitre précédent, ce chapitre a pour objectif d'identifier et de formaliser les besoins auxquels doit répondre la plateforme GREENLY. Nous y présentons successivement les besoins fonctionnels et non fonctionnels du système, le diagramme de cas d'utilisation qui en découle, ainsi que la démarche méthodologique Agile Scrum adoptée pour planifier et conduire le développement du projet.

\section{Besoins fonctionnels}

L'analyse du système a permis d'identifier les besoins fonctionnels suivants, regroupés par domaine~:

\begin{enumerate}
    \item \textbf{Gestion des comptes utilisateurs}
    \begin{itemize}
        \item[BF1.] L'utilisateur doit pouvoir créer un compte (inscription par courriel/mot de passe).
        \item[BF2.] L'utilisateur doit pouvoir se connecter et se déconnecter de manière sécurisée.
        \item[BF3.] L'utilisateur doit pouvoir réinitialiser son mot de passe en cas d'oubli.
        \item[BF4.] L'utilisateur doit pouvoir consulter et modifier les informations de son profil.
    \end{itemize}
    \item \textbf{Calcul de l'empreinte environnementale}
    \begin{itemize}
        \item[BF5.] L'utilisateur doit pouvoir saisir des données d'usage pour six catégories~: transport, énergie, électronique, eau, alimentation et déchets.
        \item[BF6.] Le système doit valider les données saisies (bornes min/max, types) avant tout calcul.
        \item[BF7.] Le système doit renvoyer une prédiction chiffrée (kg de CO2e ou litres d'eau) obtenue par un modèle de Machine Learning entraîné, et non par une simple formule statique.
        \item[BF8.] Le système doit calculer un \textit{Green Score} (0 à 100) situant le résultat de l'utilisateur par rapport à la distribution des valeurs observées lors de l'entraînement.
        \item[BF9.] Le système doit permettre d'enregistrer chaque analyse effectuée dans l'historique de l'utilisateur.
    \end{itemize}
    \item \textbf{Suivi et historique}
    \begin{itemize}
        \item[BF10.] L'utilisateur doit pouvoir consulter l'historique de toutes ses analyses passées, filtrable par catégorie et par période.
        \item[BF11.] L'utilisateur doit pouvoir supprimer une analyse de son historique.
        \item[BF12.] L'utilisateur doit pouvoir exporter un rapport de ses analyses au format PDF.
    \end{itemize}
    \item \textbf{Tableaux de bord et Business Intelligence}
    \begin{itemize}
        \item[BF13.] Le système doit présenter un tableau de bord synthétique (indicateurs clés, tendance mensuelle, répartition par catégorie).
        \item[BF14.] Le système doit fournir un tableau de bord Business Intelligence avancé, avec filtres par période et par catégorie, comparaison annuelle, et synthèse générée automatiquement à partir des agrégats réels de l'utilisateur.
        \item[BF15.] Le système doit permettre l'export des données du tableau de bord au format CSV.
    \end{itemize}
    \item \textbf{Intelligence artificielle et prévision}
    \begin{itemize}
        \item[BF16.] Le système doit afficher, pour chaque catégorie, les performances réelles du modèle de Machine Learning utilisé (algorithme retenu, MAE, RMSE, R²).
        \item[BF17.] Le système doit générer une prévision statistique de la tendance future de l'utilisateur à partir de son historique réel, avec un intervalle de confiance.
        \item[BF18.] Le système doit fournir des recommandations personnalisées fondées sur l'importance réelle des variables (\textit{feature importance}) du modèle concerné.
    \end{itemize}
\end{enumerate}

\section{Besoins non fonctionnels}

\begin{itemize}[label=\textbullet,font=\normalsize]
    \item \textbf{Performance}~: une prédiction du modèle de Machine Learning doit être renvoyée par l'API en quelques dizaines de millisecondes, les modèles étant chargés en mémoire au démarrage du service (\textit{singleton} \texttt{ModelRegistry}) plutôt que rechargés à chaque requête.
    \item \textbf{Sécurité}~: les données de chaque utilisateur doivent être strictement cloisonnées grâce à des politiques de sécurité au niveau des lignes (\textit{Row Level Security}) appliquées sur la base de données, garantissant qu'un utilisateur ne peut accéder qu'à ses propres analyses.
    \item \textbf{Fiabilité}~: les entrées utilisateur doivent être validées côté client et côté serveur (schémas Pydantic) avant tout traitement, afin d'éviter les erreurs de calcul ou les injections de données invalides.
    \item \textbf{Disponibilité}~: les services (application web et micro-service ML) doivent exposer un point de contrôle de santé (\textit{health check}) permettant une supervision automatisée en environnement conteneurisé.
    \item \textbf{Portabilité}~: l'ensemble du système doit être déployable de façon reproductible via des conteneurs Docker, indépendamment de l'environnement hôte.
    \item \textbf{Ergonomie}~: l'interface doit être responsive (adaptable mobile/bureau), accessible en mode clair et sombre, et fournir un retour visuel immédiat (indicateurs de chargement, messages d'erreur explicites).
    \item \textbf{Transparence \& traçabilité}~: chaque facteur d'émission ou constante physique utilisée pour construire les données d'entraînement doit être documentée et rattachée à une source officielle vérifiable.
    \item \textbf{Évolutivité}~: l'architecture doit permettre l'ajout de nouvelles catégories de calcul ou de nouveaux algorithmes de Machine Learning sans remise en cause de l'architecture existante (registre de modèles générique, pipelines scikit-learn interchangeables).
\end{itemize}

\section{Diagramme de cas d'utilisation}

La figure~\ref{fig:use-case} présente le diagramme de cas d'utilisation général de la plateforme GREENLY. Deux acteurs interagissent avec le système~: le \textbf{Visiteur}, qui n'a pas encore de compte, et l'\textbf{Utilisateur authentifié}, qui dispose d'un accès complet aux fonctionnalités de calcul, de suivi et d'analyse.

\begin{figure}[H]
\centering
\resizebox{\textwidth}{!}{
\begin{tikzpicture}[>=Stealth]

\tikzset{
  usecase/.style={draw, ellipse, minimum width=3.7cm, minimum height=1.15cm, align=center, font=\scriptsize, fill=isiBlue!6, line width=0.6pt},
  actorlbl/.style={font=\small\bfseries, align=center}
}

% ----- System boundary -----
\draw[line width=0.8pt] (3,-0.5) rectangle (12.6,-14.3);
\node[font=\bfseries\small] at (7.8,-0.9) {Plateforme GREENLY};

% ----- Use cases -----
\node[usecase] (uc1)  at (5.7,-2)    {S'inscrire};
\node[usecase] (uc2)  at (10.3,-2)   {Se connecter};
\node[usecase] (uc3)  at (5.7,-4.1)  {Réinitialiser le\\mot de passe};
\node[usecase] (uc4)  at (10.3,-4.1) {Effectuer un calcul\\d'empreinte};
\node[usecase] (uc5)  at (5.7,-6.2)  {Obtenir une\\prédiction ML};
\node[usecase] (uc6)  at (10.3,-6.2) {Enregistrer\\l'analyse};
\node[usecase] (uc7)  at (5.7,-8.3)  {Consulter l'historique\\des analyses};
\node[usecase] (uc8)  at (10.3,-8.3) {Consulter le\\tableau de bord};
\node[usecase] (uc9)  at (5.7,-10.4) {Consulter le tableau\\de bord BI};
\node[usecase] (uc10) at (10.3,-10.4){Consulter analyses\\et prévisions IA};
\node[usecase] (uc11) at (5.7,-12.5) {Exporter un\\rapport PDF};
\node[usecase] (uc12) at (10.3,-12.5){Gérer profil et\\paramètres};

% ----- Actors (drawn top-down: tête, corps, bras, jambes) -----
\draw (0.9,-1.6) circle (0.28);
\draw (0.9,-1.88) -- (0.9,-2.75);
\draw (0.55,-2.2) -- (1.25,-2.2);
\draw (0.9,-2.75) -- (0.55,-3.35);
\draw (0.9,-2.75) -- (1.25,-3.35);
\coordinate (visiteurPt) at (0.9,-2.2);
\node[actorlbl] at (0.9,-3.7) {Visiteur};

\draw (0.9,-8.3) circle (0.28);
\draw (0.9,-8.58) -- (0.9,-9.45);
\draw (0.55,-8.9) -- (1.25,-8.9);
\draw (0.9,-9.45) -- (0.55,-10.05);
\draw (0.9,-9.45) -- (1.25,-10.05);
\coordinate (userPt) at (0.9,-8.9);
\node[actorlbl, align=center] at (0.9,-10.4) {Utilisateur\\authentifié};

% ----- Associations : Visiteur -----
\draw (visiteurPt) -- (uc1.west);
\draw (visiteurPt) -- (uc2.west);
\draw (visiteurPt) -- (uc3.west);

% ----- Associations : Utilisateur authentifié -----
\draw (userPt) -- (uc4.west);
\draw (userPt) -- (uc7.west);
\draw (userPt) -- (uc8.west);
\draw (userPt) -- (uc9.west);
\draw (userPt) -- (uc10.west);
\draw (userPt) -- (uc11.west);
\draw (userPt) -- (uc12.west);

% ----- Include relationships -----
\draw[dashed, ->] (uc4) -- node[above, sloped, font=\tiny] {\guillemotleft include\guillemotright} (uc5);
\draw[dashed, ->] (uc4) -- node[above, sloped, font=\tiny] {\guillemotleft include\guillemotright} (uc6);

\end{tikzpicture}
}
\caption{Diagramme de cas d'utilisation général de GREENLY}
\label{fig:use-case}
\end{figure}

L'utilisateur authentifié hérite implicitement de toutes les capacités du visiteur (un compte doit avoir été créé et une connexion établie au préalable pour accéder aux fonctionnalités réservées). Le cas d'utilisation \textit{Effectuer un calcul d'empreinte} inclut systématiquement l'obtention d'une prédiction auprès du service de Machine Learning ainsi que l'enregistrement de l'analyse dans l'historique de l'utilisateur.

\section{Méthodologie Agile Scrum}

Le développement de GREENLY a été conduit selon le cadre méthodologique \textbf{Scrum} \cite{scrumGuide}, un cadre Agile itératif et incrémental particulièrement adapté à un projet à forte composante exploratoire (choix des algorithmes de Machine Learning, itérations sur l'ergonomie des tableaux de bord). Cette méthodologie repose sur~:

\begin{itemize}[label=\textbullet,font=\normalsize]
    \item un \textbf{Product Backlog}~: une liste priorisée de toutes les fonctionnalités attendues du produit, formulées sous forme de user stories~;
    \item des \textbf{Sprints}~: des itérations de développement de durée fixe (deux semaines), à l'issue desquelles un incrément fonctionnel du produit est livré~;
    \item des cérémonies Scrum~: \textit{Sprint Planning} (planification en début de sprint), \textit{Daily Scrum} (point d'avancement quotidien), \textit{Sprint Review} (démonstration de l'incrément) et \textit{Sprint Retrospective} (amélioration continue du processus)~;
    \item des rôles définis~: le \textit{Product Owner} (responsable de la priorisation du backlog selon la valeur métier), le \textit{Scrum Master} (garant du respect du processus) et l'équipe de développement.
\end{itemize}

Dans le cadre de ce projet à développeur unique, ces rôles ont été assumés de manière combinée, avec un point de synchronisation hebdomadaire avec l'encadrant académique jouant le rôle de \textit{Product Owner} pour la validation des priorités.

\section{Product Backlog}

Le tableau~\ref{tab:backlog} présente un extrait représentatif du Product Backlog du projet, avec pour chaque item sa priorité selon la méthode MoSCoW (\textit{Must have}, \textit{Should have}, \textit{Could have}) et son estimation en points d'effort (échelle de Fibonacci~: 1, 2, 3, 5, 8, 13).

\begin{table}[H]
\centering
\caption{Extrait du Product Backlog}
\label{tab:backlog}
\begin{tabular}{|p{0.7cm}|p{7.6cm}|p{1.9cm}|p{1.3cm}|}
\hline
\rowcolor{isiBlue!15}
\textbf{ID} & \textbf{User Story} & \textbf{Priorité} & \textbf{Effort} \\
\hline
US1  & En tant que visiteur, je veux créer un compte afin d'accéder à la plateforme. & Must have & 3 \\ \hline
US2  & En tant qu'utilisateur, je veux me connecter de manière sécurisée. & Must have & 3 \\ \hline
US3  & En tant qu'utilisateur, je veux réinitialiser mon mot de passe en cas d'oubli. & Should have & 2 \\ \hline
US4  & En tant que Data Scientist, je veux collecter et documenter des jeux de données publics réels pour chaque catégorie. & Must have & 8 \\ \hline
US5  & En tant que Data Scientist, je veux prétraiter les données (nettoyage, encodage) pour les rendre exploitables. & Must have & 5 \\ \hline
US6  & En tant que Data Scientist, je veux comparer plusieurs algorithmes de régression et retenir automatiquement le meilleur par catégorie. & Must have & 13 \\ \hline
US7  & En tant qu'utilisateur, je veux saisir mes données de trajet et obtenir mon empreinte carbone transport. & Must have & 8 \\ \hline
US8  & En tant qu'utilisateur, je veux effectuer le même calcul pour l'énergie, l'électronique, l'eau, l'alimentation et les déchets. & Must have & 13 \\ \hline
US9  & En tant qu'utilisateur, je veux voir un Green Score qui situe mon résultat par rapport aux autres. & Should have & 5 \\ \hline
US10 & En tant qu'utilisateur, je veux consulter l'historique de mes analyses passées. & Must have & 5 \\ \hline
US11 & En tant qu'utilisateur, je veux un tableau de bord résumant mon impact global. & Must have & 8 \\ \hline
US12 & En tant qu'utilisateur, je veux un tableau de bord BI avec filtres, comparaison annuelle et synthèse automatique. & Should have & 13 \\ \hline
US13 & En tant qu'utilisateur, je veux visualiser une prévision de tendance basée sur mon historique réel. & Should have & 8 \\ \hline
US14 & En tant qu'utilisateur, je veux recevoir des recommandations personnalisées fondées sur mes données. & Should have & 5 \\ \hline
US15 & En tant qu'utilisateur, je veux exporter mes analyses en PDF. & Could have & 3 \\ \hline
US16 & En tant qu'administrateur système, je veux déployer l'application via Docker pour garantir la reproductibilité. & Must have & 5 \\ \hline
\end{tabular}
\end{table}

\section{Sprints}

Le Product Backlog a été découpé en six sprints de deux semaines, précédés d'un Sprint~0 de cadrage technique. Le tableau~\ref{tab:sprints} résume le contenu de chaque sprint.

\begin{table}[H]
\centering
\caption{Répartition du backlog par sprint}
\label{tab:sprints}
\begin{tabular}{|p{1.6cm}|p{9.3cm}|p{2.3cm}|}
\hline
\rowcolor{isiBlue!15}
\textbf{Sprint} & \textbf{Objectif et user stories couvertes} & \textbf{Durée} \\
\hline
Sprint 0 & Cadrage technique~: choix de la stack (React/TypeScript, FastAPI, Supabase), mise en place des dépôts et de l'environnement Docker. & 1 semaine \\ \hline
Sprint 1 & Authentification et modélisation de la base de données (US1, US2, US3) -- schéma \texttt{analyses}, politiques RLS. & 2 semaines \\ \hline
Sprint 2 & Collecte et préparation des données réelles, premier pipeline d'entraînement pour la catégorie Transport (US4, US5, US6 partiel, US7). & 2 semaines \\ \hline
Sprint 3 & Extension du pipeline aux cinq autres catégories, comparaison complète des 5 algorithmes, intégration de l'API de prédiction (US6, US8). & 2 semaines \\ \hline
Sprint 4 & Tableau de bord principal, historique des analyses, export PDF (US9, US10, US11, US15). & 2 semaines \\ \hline
Sprint 5 & Tableau de bord Business Intelligence, module de prévision, recommandations fondées sur le \textit{feature importance} (US12, US13, US14). & 2 semaines \\ \hline
Sprint 6 & Conteneurisation Docker, tests de bout en bout, corrections et polissage de l'interface (US16). & 2 semaines \\ \hline
\end{tabular}
\end{table}

\section{Planning}

La figure~\ref{fig:gantt} présente le planning prévisionnel du projet sous forme de diagramme de Gantt, couvrant treize semaines de développement.

\begin{figure}[H]
\centering
\resizebox{\textwidth}{!}{
\begin{tikzpicture}[x=0.95cm, y=0.9cm]

% Grid / week axis
\foreach \w in {0,...,13} {
    \draw[gray!35] (\w,0) -- (\w,-7.5);
}
\foreach \w in {1,...,13} {
    \node[font=\tiny] at (\w-0.5,0.4) {S\w};
}
\draw[line width=0.6pt] (0,0) -- (13,0);

% Task rows: nom du sprint / semaine de début / durée (semaines) / numéro de ligne
\foreach \name/\start/\dur/\row in {
    Sprint 0 -- Cadrage technique/0/1/1,
    Sprint 1 -- Auth \& base de données/1/2/2,
    Sprint 2 -- Data \& 1er pipeline ML/3/2/3,
    Sprint 3 -- Extension ML (6 catégories)/5/2/4,
    Sprint 4 -- Dashboard historique et export PDF/7/2/5,
    Sprint 5 -- BI prévision et recommandations/9/2/6,
    Sprint 6 -- Déploiement \& tests/11/2/7
} {
    \pgfmathsetmacro{\ypos}{-\row*0.95}
    \draw[fill=isiBlue!55, draw=isiBlue, line width=0.4pt]
        (\start,\ypos) rectangle ++(\dur,0.6);
    \node[font=\tiny, anchor=west] at (13.2,\ypos+0.3) {\name};
}

\end{tikzpicture}
}
\caption{Diagramme de Gantt du projet GREENLY}
\label{fig:gantt}
\end{figure}

\section*{Conclusion}
Ce chapitre a permis de spécifier précisément les besoins fonctionnels et non fonctionnels de la plateforme GREENLY, de les traduire en un diagramme de cas d'utilisation, puis de les organiser au sein d'un Product Backlog planifié sur six sprints selon la méthodologie Agile Scrum. Cette base de spécification constitue le point de départ de la phase de conception, détaillée dans le chapitre suivant pour le volet Intelligence Artificielle.
